\documentclass[10pt,a4paper]{amsart}
\usepackage[margin=19mm]{geometry}
\usepackage{amsmath,amssymb,mathtools}
\usepackage[scr=rsfs,cal=euler]{mathalfa}
\usepackage{xfrac}
\usepackage[unicode,hidelinks,bookmarks=false]{hyperref}
\newcommand{\ie}{i.e.\ }
\newcommand{\cf}{cf.\ }
\newcommand{\resp}{respectively\ }
\newcommand{\Subsection}{\subsection}
\providecommand{\nobreakdash}{\nobreak\hbox{-}}
\setlength{\emergencystretch}{2em}
\multlinegap0pt
\usepackage{bm}
\usepackage{bbm}
\usepackage{dsfont}
\usepackage{xspace}
\usepackage{cancel}
\usepackage{mathtools}
\usepackage{amsthm}
\makeatletter
\@ifundefined{theorem}{\newtheorem{theorem}{Theorem}}{}
\@ifundefined{lemma}{\newtheorem{lemma}[theorem]{Lemma}}{}
\makeatother
\title[Simplicial Homology Groups]{Simplicial Homology Groups}
\author{Sanjay Mishra}
\date{Source: arXiv:2511.03326v1 (CC BY 4.0)}
\begin{document}
\maketitle

\noindent\emph{Source and licence.} Simplicial Homology Groups. Version arXiv:2511.03326v1 (CC BY 4.0), distributed under the Creative Commons Attribution 4.0 International licence, as recorded in the arXiv licence metadata for this exact version and shown on the abstract page https://arxiv.org/abs/2511.03326v1. Full rights notice: licenses/arxiv-2511.03326-CC-BY-4.0.txt.\par\smallskip

\noindent\textit{Editorial scope.} Counted unit: the Lemma whose source label is \texttt{l:Computing Bases of Chain Groups via Elementary Chains} in the article (source0.tex lines 849--851, source label \texttt{l:Computing Bases of Chain Groups via Elementary Chains}), delivered together with the article's own complete original proof of it (source0.tex lines 852--916, one \texttt{proof} environment of 65 lines). No mathematics is written, repaired or re-derived: statement and proof are the article's own text line for line. Required context retained uncounted: none. Every definition, display and label of those lines is retained unchanged. Omitted from this package and not redistributed: 1--848, 917--2105. Nothing inside a retained excerpt is omitted or altered: every retained body is delivered exactly as the article's lines are written, so no omission receipt of any kind accompanies this package. Citations: the retained excerpts carry no citation command at all, so no bibliography entry of the article is transcribed and no reference is redistributed. Reference adaptation: none was needed and none was made; every reference inside the delivered unit resolves inside it, so no unresolved reference or citation is produced. Printed slips kept literal: no printed word, symbol, hypothesis or number of the counted statement or of its proof was altered. Layout and class: the article's own class is \texttt{amsart} with packages including array, amssymb, tikz, amsmath, bm, color, graphicx, hyperref, anysize, multicol, mathtools, enumitem, natbib, comment; the wrapper is the lane's pinned \texttt{amsart} editorial wrapper at 10pt on A4, which restates the theorem-like environments the retained text needs and adds the article's own macro definitions that the retained text needs (none), with extra wrapper packages (bm, bbm, dsfont, xspace, cancel, mathtools). The delivered numbering is the unit's own. Assets: no retained span contains a figure environment, a graphics-inclusion command, a PDF-page inclusion, a table or a TikZ picture; no image file is copied into this package, the package declares no asset dependency and no image file is redistributed with it. Licence: version arXiv:2511.03326v1 of this article is distributed under the Creative Commons Attribution 4.0 International licence, as recorded in the arXiv licence metadata for this exact version and shown on the abstract page https://arxiv.org/abs/2511.03326v1. A full rights notice is delivered at licenses/arxiv-2511.03326-CC-BY-4.0.txt. Characters: every retained excerpt is pure ASCII, so the delivered engine, XeLaTeX with its default Latin Modern text font, reports no missing character.
\begin{lemma}[Computing Bases of Chain Groups via Elementary Chains]\label{l:Computing Bases of Chain Groups via Elementary Chains}
A basis for $p$-chain group $C_{p}(K)$ can be obtained by orienting each $p$-simplex and using the corresponding elementary chains as a basis.   
\end{lemma}

\begin{proof}
Let $K$ be a simplicial complex, and let $\widetilde{K}_p$ denote the set of oriented $p$-simplices of $K$. The goal of this proof is to show that the set of elementary chains corresponding to the oriented $p$-simplices forms a basis for the $p$-chain group $C_p(K)$.

\textbf{Step 1: Definition of $p$-chain group $C_p(K)$}

The \textit{$p$-chain group} $C_p(K)$ consists of formal sums of oriented $p$-simplices with integer coefficients. That is, an element $c \in C_p(K)$ can be written as:
\[
c = \sum_{\sigma \in \widetilde{K}_p} a_\sigma \cdot \sigma,
\]
where:
- $\sigma$ is an oriented $p$-simplex from the set $\widetilde{K}_p$,
- $a_\sigma \in \mathbb{Z}$ are integer coefficients.

Thus, $C_p(K)$ is the free abelian group generated by the set of oriented $p$-simplices $\widetilde{K}_p$, and each element of $C_p(K)$ is a formal sum (linear combination) of these oriented simplices.

\textbf{Step 2: Definition of Elementary Chains}

For each oriented $p$-simplex $\sigma \in \widetilde{K}_p$, we define the \textit{elementary chain} $\delta_\sigma$ as a function from the set of all oriented simplices to the integers $\mathbb{Z}$, which is defined as:
\[
\delta_\sigma(\tau) = 
\begin{cases}
1 & \text{if } \tau = \sigma, \\
-1 & \text{if } \tau = \bar{\sigma}, \text{ where } \bar{\sigma} \text{ is the opposite orientation of } \sigma, \\
0 & \text{otherwise}.
\end{cases}
\]
Thus, $\delta_\sigma$ is a function that "picks out" the oriented simplex $\sigma$ and assigns it the coefficient 1, assigns the opposite orientation $\bar{\sigma}$ the coefficient -1, and assigns all other simplices the coefficient 0.

\textbf{Step 3: The Set of Elementary Chains Forms a Free Abelian Group}

The set of elementary chains $\{ \delta_\sigma : \sigma \in \widetilde{K}_p \}$ forms a \textit{free abelian group}. To show this, we must verify the following two conditions:
1. \textit{Linear Independence:}  
   If a linear combination of elementary chains is equal to the zero chain, then the coefficients must all be zero. More formally, if:
   \[
   \sum_{i} a_i \delta_{\sigma_i} = 0,
   \]
   then $a_i = 0$ for all $i$. This is because the elementary chains are defined to be nonzero only on the specific simplices they correspond to, so for the sum to be zero, each coefficient must vanish.
   
2. \textit{Free Group Structure:}  
   There is a unique way to express any $p$-chain as a linear combination of elementary chains. Given that the elementary chains are linearly independent, they generate a free abelian group.

Thus, the set $\{ \delta_\sigma : \sigma \in \widetilde{K}_p \}$ is a \textit{basis} for the free abelian group $C_p(K)$.

\textbf{Step 4: Orientation of $p$-simplices}

Each $p$-simplex $\sigma \in K_p$ is \textit{oriented} in the sense that it is given a direction, which leads to an ordered set of vertices. The orientation distinguishes $\sigma$ from its opposite $\bar{\sigma}$. Since we consider the elementary chain associated with each \textit{oriented} $p$-simplex, the basis elements $\delta_\sigma$ correspond to these oriented simplices.

The fact that the simplices are oriented ensures that each simplex has a distinct representation in the chain group, as $\sigma$ and $\bar{\sigma}$ are treated as different elements in the basis.

\textbf{Step 5: Linear Independence and Spanning}

The set $\{ \delta_\sigma : \sigma \in \widetilde{K}_p \}$ is linearly independent, as shown in Step 3. Furthermore, any $p$-chain $c \in C_p(K)$ can be expressed as a unique linear combination of these elementary chains. Specifically, any $p$-chain $c$ can be written as:
\[
c = \sum_{\sigma \in \widetilde{K}_p} a_\sigma \cdot \delta_\sigma,
\]
where the coefficients $a_\sigma \in \mathbb{Z}$ correspond to the integer coefficients of the simplices in the chain.

Thus, the set $\{ \delta_\sigma : \sigma \in \widetilde{K}_p \}$ \textit{spans} the $p$-chain group $C_p(K)$.

\textbf{Step 6: Conclusion}

Since the set $\{ \delta_\sigma : \sigma \in \widetilde{K}_p \}$ is both linearly independent and spans $C_p(K)$, it forms a basis for the $p$-chain group. Therefore, we have shown that:

A basis for  $C_p(K)$ can be obtained by orienting each $p$-simplex and using the corresponding elementary chains as a basis.
\end{proof}

\end{document}
